\section*{Conclusion}
\addcontentsline{toc}{section}{Conclusion}
\label{sec:conclusion}

AI responsibility does not disappear when no person selects the harmful output. It moves into the deployment. Organizations choose objectives, models, interfaces, memory, data, tools, user populations, evaluations, monitoring, and stop conditions. Those choices are observable, governable, and legally consequential even when the model's internal path is difficult to explain.

The deployment-control cascade makes that responsibility visible. It traces authority and resources toward operation, feedback toward those able to respond, and override or shutdown power to the point where intervention can work. Its process model, control actions, feedback, and records translate into the facts ordinary law already uses: authorization, notice, foreseeability, feasible precaution, causation, and proof. Probabilistic output changes the location of the evidence; it does not create a responsibility vacuum.

The two-layer architecture turns that map into an institutional rule. For designated high-impact deployments, an enterprise with material control over the risk should face the claimant under a nondelegable deployment duty. In a divided stack, the identification rule reveals which enterprises governed which risk. In every covered deployment, the production rule supplies the proportionate record needed to test the merits. The claimant still proves recognized harm, breach or defect, causation, and every other substantive element. The front door makes those rights provable without converting them into strict liability.

Inside the enterprise, the responsibility charter aligns accountability with authority before harm occurs. It names risk owners, fixes approval, records dissent, protects escalation, tests stop and rollback power, and preserves incident evidence. The allocation protects the person who reports a danger and preserves recourse against the person who conceals, falsifies, retaliates, or bypasses a constraint. Public responsibility and internal fairness reinforce the same information system.

The anti-substitution principle carries that logic into care, professional judgment, accommodation, and public service. An automatic instrumentality may perform a legally required function and often improve it. Its mere availability does not perform the duty. The answerable institution must preserve the protection the duty exists to secure, including an effective path to notice, review, intervention, or correction where the governing law requires one.

None of this requires a machine defendant. Unless positive law constructs a new juridical platform, the legal subjects remain the natural persons and organizations capable of holding rights, bearing duties, producing records, and answering remedies. Functional classifications may still vary: the same deployed system can be a product, service, instrument, or communicative medium for different rules. The stable point is responsibility. Courts should classify the function, reconstruct the control path, open the evidentiary front door, and keep the obligation attached to an actor law can reach.
